\documentclass[10pt,a4paper]{amsart}
\usepackage{amsmath,amssymb,amsthm,graphicx}
\usepackage{verbatim}
\newenvironment{reference}{\comment}{\endcomment}
%\newenvironment{reference}{}{}
\newenvironment{slogan}{\comment}{\endcomment}
\newenvironment{history}{\comment}{\endcomment}

% For commutative diagrams we use Xy-pic
\usepackage[all]{xy}

% We use 2cell for 2-commutative diagrams.
\xyoption{2cell}
% Theorem environments.
%
\theoremstyle{plain}
\newtheorem{theorem}[subsection]{Theorem}
\newtheorem{proposition}[subsection]{Proposition}
\newtheorem{lemma}[subsection]{Lemma}

\theoremstyle{definition}
\newtheorem{definition}[subsection]{Definition}
\newtheorem{example}[subsection]{Example}
\newtheorem{exercise}[subsection]{Exercise}
\newtheorem{situation}[subsection]{Situation}

\theoremstyle{remark}
\newtheorem{remark}[subsection]{Remark}
\newtheorem{remarks}[subsection]{Remarks}

\numberwithin{equation}{subsection}

% Macros
%
\def\lim{\mathop{\mathrm{lim}}\nolimits}
\def\colim{\mathop{\mathrm{colim}}\nolimits}
\def\Spec{\mathop{\mathrm{Spec}}}
\def\Hom{\mathop{\mathrm{Hom}}\nolimits}
\def\Ext{\mathop{\mathrm{Ext}}\nolimits}
\def\SheafHom{\mathop{\mathcal{H}\!\mathit{om}}\nolimits}
\def\SheafExt{\mathop{\mathcal{E}\!\mathit{xt}}\nolimits}
\def\Sch{\mathit{Sch}}
\def\Mor{\mathop{\mathrm{Mor}}\nolimits}
\def\Ob{\mathop{\mathrm{Ob}}\nolimits}
\def\Sh{\mathop{\mathit{Sh}}\nolimits}
\def\NL{\mathop{N\!L}\nolimits}
\def\CH{\mathop{\mathrm{CH}}\nolimits}
\def\proetale{{pro\text{-}\acute{e}tale}}
\def\etale{{\acute{e}tale}}
\def\QCoh{\mathit{QCoh}}
\def\Ker{\mathop{\mathrm{Ker}}}
\def\Im{\mathop{\mathrm{Im}}}
\def\Coker{\mathop{\mathrm{Coker}}}
\def\Coim{\mathop{\mathrm{Coim}}}

% Boxtimes
%
\DeclareMathSymbol{\boxtimes}{\mathbin}{AMSa}{"02}

%
% Macros for moduli stacks/spaces
%
\def\QCohstack{\mathcal{QC}\!\mathit{oh}}
\def\Cohstack{\mathcal{C}\!\mathit{oh}}
\def\Spacesstack{\mathcal{S}\!\mathit{paces}}
\def\Quotfunctor{\mathrm{Quot}}
\def\Hilbfunctor{\mathrm{Hilb}}
\def\Curvesstack{\mathcal{C}\!\mathit{urves}}
\def\Polarizedstack{\mathcal{P}\!\mathit{olarized}}
\def\Complexesstack{\mathcal{C}\!\mathit{omplexes}}
% \Pic is the operator that assigns to X its picard group, usage \Pic(X)
% \Picardstack_{X/B} denotes the Picard stack of X over B
% \Picardfunctor_{X/B} denotes the Picard functor of X over B
\def\Pic{\mathop{\mathrm{Pic}}\nolimits}
\def\Picardstack{\mathcal{P}\!\mathit{ic}}
\def\Picardfunctor{\mathrm{Pic}}
\def\Deformationcategory{\mathcal{D}\!\mathit{ef}}
\usepackage{cleveref}
\crefname{originalreference}{original reference}{original references}
\Crefname{originalreference}{Original reference}{Original references}
\makeatletter
\newlabel{lemma-w-contractible-is-weakly-contractible}{{098H}{source}{}{sourcepointer-1}{}}
\newlabel{lemma-w-contractible-is-weakly-contractible@cref}{{[originalreference][0][]098H}{[1][0][]source}{}{}{}}
\makeatother
\crefname{originaltheorem}{original theorem}{original theorems}
\Crefname{originaltheorem}{Original theorem}{Original theorems}
\makeatletter
\newlabel{sites-theorem-plus}{{00WB}{source}{}{sourcepointer-2}{}}
\newlabel{sites-theorem-plus@cref}{{[originaltheorem][0][]00WB}{[1][0][]source}{}{}{}}
\makeatother
\title{A presheaf sending finite disjoint unions to products agrees with its sheafification on affine weakly contractible objects}
\author{The Stacks Project Authors}
\date{stacks-project:proetale.tex:local-snapshot-20261012; GNU FDL 1.2 or later}
\begin{document}
\maketitle
\noindent\textit{Editorial note.} All delivered mathematical text is the author's own, copied from the pinned Stacks Project chapter file proetale.tex; no statement or proof step was rewritten, reordered or removed.\par\medskip
\noindent\textit{Editorial note.} The counted proof is the author's complete argument: the reduction to the plus construction, the verification that the plus construction still sends finite disjoint unions to products, and the computation of the value on an affine weakly contractible object through the refinement of its pro-etale coverings.\par\medskip
\noindent\textit{Editorial note.} The label lemma-presheaf-value-weakly-contractible names the counted statement, and the proof environment immediately following it in the source is delivered verbatim as its proof.\par\medskip
\noindent\textit{Editorial note.} The counted statement is carried by the author's own lemma environment in the chapter source.\par\medskip
\noindent\textit{Editorial note.} Pre-open gate: python operations/agent-authoring/tools/precard\_gate.py --source proetale.tex --label lemma-presheaf-value-weakly-contractible --start 3464 --end 3497 --local-id 3591 returned RESULT=OK (L3591 free, occupancy FREE, public interval clear). Candidate taken from sources/stacks-proetale/candidates-halfA.json.\par\medskip
\noindent\textit{Editorial note.} The statement and proof are copied from the pinned chapter file; the cited inputs are kept as external pointers with their canonical Stacks tags.\par\medskip
\noindent\textit{Allowed standard prerequisites (editorial).}\begin{itemize}
\item Description of the sheafification as a double plus construction. Sites, the theorem identifying the sheafification of a presheaf with the twice applied plus construction. Supplies the reduction to proving that the plus construction preserves the hypotheses and the value on an affine weakly contractible object.
\item Refinement of pro-etale coverings of a weakly contractible object. Same chapter, the lemma that every pro-etale covering of an object satisfying the refinement property can be refined by a finite Zariski covering of the form a disjoint union. Supplies the finite disjoint union decomposition refining any pro-etale covering of W, which is the key input for the second part.
\end{itemize}\medskip
\ifdefined\hypertarget\hypertarget{sourcepointer-1}{}\fi
\noindent\textit{Source reference.} Original reference 098H: sources/stacks-proetale/pinned-original-source.tex, source lines 3027--3027. Prerequisite pointer invoked at proof line(s) 3493; it names a standard fact of the same chapter that the counted argument uses, and every step of the counted proof built on it is retained verbatim in the delivered body.\par\medskip
\ifdefined\hypertarget\hypertarget{sourcepointer-2}{}\fi
\noindent\textit{Source reference.} Original theorem 00WB: sources/stacks-chow/tags.txt, source lines 1177--1177. Prerequisite pointer invoked at proof line(s) 3474; it names a standard fact of a sibling chapter that the counted argument uses, and every step of the counted proof built on it is retained verbatim in the delivered body.\par\medskip
\begin{lemma}
\label{lemma-presheaf-value-weakly-contractible}
Let $X$ be a scheme. Let $\mathcal{F}$ be a presheaf of sets on $X_\proetale$
which sends finite disjoint unions to products. Then
$\mathcal{F}^\#(W) = \mathcal{F}(W)$ if $W$ is an affine weakly contractible
object of $X_\proetale$.
\end{lemma}

\begin{proof}
Recall that $\mathcal{F}^\#$ is equal to $(\mathcal{F}^+)^+$, see
Sites, Theorem \ref{sites-theorem-plus}, where $\mathcal{F}^+$
is the presheaf which sends an object $U$ of $X_\proetale$ to
$\colim H^0(\mathcal{U}, \mathcal{F})$ where the colimit is
over all pro-\'etale coverings $\mathcal{U}$ of $U$.
Thus it suffices to prove that (a) $\mathcal{F}^+$ sends finite
disjoint unions to products and (b) sends $W$ to $\mathcal{F}(W)$.
If $U = U_1 \amalg U_2$, then given a pro-\'etale covering
$\mathcal{U} = \{f_j : V_j \to U\}$ of $U$ we obtain pro-\'etale coverings
$\mathcal{U}_i = \{f_j^{-1}(U_i) \to U_i\}$ and we clearly have
$$
H^0(\mathcal{U}, \mathcal{F}) =
H^0(\mathcal{U}_1, \mathcal{F}) \times
H^0(\mathcal{U}_2, \mathcal{F})
$$
because $\mathcal{F}$ sends finite disjoint unions to products (this includes
the condition that $\mathcal{F}$ sends the empty scheme to the singleton).
This proves (a).
Finally, any pro-\'etale covering of $W$ can be refined by a finite disjoint
union decomposition $W = W_1 \amalg \ldots W_n$ by
Lemma \ref{lemma-w-contractible-is-weakly-contractible}.
Hence $\mathcal{F}^+(W) = \mathcal{F}(W)$ exactly because the value of
$\mathcal{F}$ on $W$ is the product of the values of $\mathcal{F}$
on the $W_j$. This proves (b).
\end{proof}

\end{document}
